\begin{frame}{Now change the question: how does it move?}
\princ{FROM VALUES TO STATES}
A \textbf{finite-state machine} is a bounded model with named states and input-triggered transitions.
\begin{tabular}{@{}p{3.0cm}p{9.5cm}@{}}
\textbf{input} & An event supplied to the machine.\\[2pt]
\textbf{start state} & The state where a trace begins.\\[2pt]
\textbf{transition} & The declared next-state rule for a state and input.\\[2pt]
\textbf{trace} & The ordered states visited as inputs are processed.\\[2pt]
\textbf{model limit} & A relevant fact the model does not represent or establish.
\end{tabular}
\predict{What must be named before another person can reproduce a trace?}
\see{states, input alphabet, start state, transitions, and input sequence.}
\end{frame}

\begin{frame}{Synthetic badge model}
\princ{MAKE TRANSITIONS EXPLICIT}
\begin{columns}[T]
\begin{column}{0.55\textwidth}
\begin{tikzpicture}[node distance=24mm, every node/.style={font=\small}, state/.style={draw,rounded corners,fill=blue!10,minimum width=2cm,minimum height=1cm}, wire/.style={-{Stealth[length=5pt]},thick}]
\node[state] (locked) {Locked};
\node[state,right=of locked] (open) {Open};
\draw[wire] (locked) to[bend left=18] node[above] {valid} (open);
\draw[wire] (open) to[bend left=18] node[below] {timeout} (locked);
\end{tikzpicture}
\end{column}
\begin{column}{0.4\textwidth}
\textbf{Declared inputs:} valid, timeout.\\[4pt]
\textbf{Start:} Locked.\\[4pt]
\textbf{Classroom rule:} other inputs leave the state unchanged.
\end{column}
\end{columns}
\predict{Starting at Locked, what states follow valid, timeout, valid?}
\see{Locked $\to$ Open $\to$ Locked $\to$ Open.}
\mistake{An input is not a state; it is the event used to choose a transition.}
\end{frame}

\begin{frame}{Trace it one input at a time}
\princ{CHECK THE TRACE}
\begin{center}
\begin{tabular}{c|c|c}
\textbf{step} & \textbf{input} & \textbf{state after input}\\\midrule
start & --- & Locked\\
1 & valid & Open\\
2 & timeout & Locked\\
3 & valid & Open
\end{tabular}
\end{center}
\begin{block}{Trace}
\code{Locked -> Open -> Locked -> Open}
\end{block}
\see{the trace includes the start state so another person can reproduce it.}
\end{frame}

\begin{frame}[fragile]{Prompt 04: run the complete model}
\princ{REPRODUCIBLE TRACE}
\promptbox{04}{a}
\predict{What should the final state be after the three inputs?}
\see{the exact four-state sequence is Locked, Open, Locked, Open.}
\end{frame}

\begin{frame}{A model can be incomplete}
\princ{FIND WHAT WAS FORGOTTEN}
\begin{columns}[T]
\begin{column}{0.48\textwidth}
\textbf{Declared states:} Locked, Open\\
\textbf{Inputs:} valid, timeout\\[6pt]
\begin{tabular}{@{}l l@{}}
\textbf{Present} & Locked + valid $\to$ Open\\
\textbf{Missing} & Open + timeout $\to$ Locked
\end{tabular}
\end{column}
\begin{column}{0.47\textwidth}
\begin{alertblock}{Do not fill it silently}
An implicit self-loop is not the same as a declared transition. List the state/input pairs and compare them with the table.
\end{alertblock}
\end{column}
\end{columns}
\predict{Which pair must a validator report?}
\see{exactly $(Open, timeout)$ is absent from the intentionally incomplete model.}
\end{frame}

\begin{frame}[fragile]{Prompt 05: name the missing transition}
\princ{VALIDATE THE MODEL}
\promptbox{05}{a}
\predict{What should appear in the missing-transition list?}
\see{exactly $(Open, timeout)$, with no automatic self-loop.}
\end{frame}
